% =========================================================
% CHAPTER 4 — EXPERIMENTAL METHODOLOGY
% =========================================================

\chapter{Methodology}
\label{chap:experimental_methodology}

The exploratory data analysis presented in Chapter~\ref{chap:eda} revealed a consistent imbalance across the three datasets under study. In all cases, the normoglycemic range represented the majority of the available measurements, while hypoglycemia, and particularly severe hypoglycemia, appeared only in a small fraction of the records. Severe hyperglycemia (>250 mg/dL) was also underrepresented compared to the normoglycemic range, although to a lesser extent. This imbalance was observed both at the global dataset level and at the patient level, suggesting that predictive models trained on these data may be biased towards the most frequently observed glycemic states. This chapter defines the experimental methodology used to address the glycemic-range imbalance identified in Chapter~\ref{chap:eda}

\section{Experimental Design}
\label{sec:experimental_methodology_proposal}

\subsection{Scope of the Balancing Process}
\label{subsec:experimental_methodology_scope}

First, it is important to note that, although the imbalance has been identified in the original CGM individual measurements, the balancing process will not be applied directly to these raw timestamped records. These measurements form actual patient-specific \textbf{time series} in which each glucose value is linked to previous and subsequent observations through a temporal dependency structure. Inserting, deleting, or duplicating individual records into this structure would create artificial measurements that do not correspond to real physiological glycemic profiles, and introduce strong biases unrelated to the actual imbalance problem being studied. For this reason, balancing will be applied only after the original timestamped measurements have been transformed into the supervised \textbf{data windows} used by the predictive models. At this stage, each instance is represented as an input-output pair derived from the original CGM sequence, which makes it possible to apply resampling techniques without directly modifying the underlying timestamped measurements structure. This preserves the temporal consistency of the raw data while allowing the learning algorithm to be exposed to a less imbalanced distribution of glycemic ranges during model fitting.

Additionally, \textbf{balancing is applied only to the training data} because the purpose of these techniques is to modify the distribution of examples seen by the model during training, so that underrepresented glycemic ranges have a greater influence on parameter estimation. Therefore, validation and test sets must remain unchanged, since they are intended to preserve the original data distribution in order to provide an unbiased estimate of how the model behaves under realistic conditions. This decision is based on the standard methodological practice found in the imbalanced learning literature explored in Chapter~\ref{chap:SoA}. Throughout the state of the art reviewed in said chapter, resampling methods are consistently applied to the training data, while validation and test subsets are kept unmodified. Also, an additional specific review of methodological recommendations for this decision did not identify evidence suggesting that validation or test sets should be balanced. On the contrary, official documentation \cite{imbalanced_learn} explicitly warns that resampling the whole dataset before splitting, or equivalently resampling both training and test partitions, leads to data leakage and to evaluation on a distribution that no longer reflects the real use case. Similarly, the scikit-learn documentation \cite{scikit_learn_common_pitfalls} states that preprocessing operations must be learned only from the training data, and that test data should not be involved in fitting procedures.

\subsection{Fixed Experimental Configuration}
\label{subsec:experimental_methodology_fixed_configuration}

In order to isolate the effect of the balancing strategies, several experimental conditions will remain fixed throughout all experiments in order to provide a stable baseline for comparison. These decisions have been chosen because they represent common standard settings in glucose prediction studies, as studied in Chapter~\ref{chap:SoA} and because the objective of this thesis is not to compare different predictive architectures, prediction horizons, or window configurations, but to evaluate how modifying the training-data distribution affects glucose prediction performance under equivalent experimental conditions.

For this reason, all experiments will use the same predictive algorithm, namely a \textbf{Long Short-Term Memory} network. LSTM models are appropriate for this task because they are designed to process sequential data and can exploit temporal dependencies between previous CGM measurements.

The prediction horizon will be fixed at \textbf{one hour} into the future. Since the supervised windows are generated from 15-minute CGM samples, this corresponds to \textbf{predicting four samples ahead}. The input history length will be fixed at \textbf{eight previous measurements}, equivalent to \textbf{two hours of past glucose information}.

Model training is performed using five-fold cross-validation. 

Finally, \textbf{Root Mean Squared Error} will be used as the main loss and evaluation metric.

\section{Data-Balancing Methodology}
\label{sec:data-balancing methodology}

\subsection{Hybrid Balancing Strategy}
\label{subsec:experimental_methodology_hybrid_strategy}

To address the numerical imbalance limitation identified in Chapter~\ref{chap:eda}, this study proposes a hybrid balancing strategy that combines synthetic oversampling of minority glycemic ranges, with controlled undersampling of the majority ranges. This approach preserves the total number of measurements in each dataset while redistributing them across glycemic ranges to achieve the desired proportions. This strategy will expose predictive models to a more representative set of clinically relevant glycemic events while preserving the overall characteristics of the original datasets.

\subsection{Target Distributions}
\label{subsec:target_distributions}

The exact proportions in which the hybrid approach described previously will change each dataset distribution are explored through three different target distributions: normalized baseline balancing, semi-balancing, and full balancing.

\subsubsection{Normalized baseline}
\label{subsubsec:normalized_baseline}

First, a normalized baseline target distribution is defined from the EDA results, providing a \textbf{single common reference} distribution across the three cohorts. This approach computes the target samples for each of the glycemic ranges as a sample-weighted mean of the three datasets percentages for that range, where each dataset contributes in proportion to its number of valid glucose samples. This decision avoids giving the same weight to datasets of very different size, as larger datasets should have a stronger influence on the common reference distribution.

This normalized baseline is obtained by pooling the glycemic-range counts of DiaTrend, REPLACE-BG, and T1DiabetesGranada, and then computing the global proportion of each range over the combined total.

Formally, let $n(d,r)$ denote the number of samples belonging to glycemic range $r$ in dataset $d$, and let $N(d)$ denote the total number of valid glucose samples in dataset $d$:

\[
N(d) = \sum_{r} n(d,r).
\]

The normalized baseline proportion for glycemic range $r$ is then defined as:

\[
p_{\mathrm{norm}}(r)
=
\frac{\displaystyle\sum_{d \in \mathcal{D}} n(d,r)}
{\displaystyle\sum_{d \in \mathcal{D}} N(d)},
\]

where $\mathcal{D}$ contains DiaTrend, REPLACE-BG, and T1DiabetesGranada.

As a result, the normalized baseline target proportions are described in Table~\ref{tab:normalized_baseline_proportions}

\begin{table}[htbp]
    \centering
    \caption{Normalized baseline target proportions.}
    \label{tab:normalized_baseline_proportions}
    \begin{tabular}{lr}
        \hline
        \textbf{Glycemic range} & \textbf{Percentage} \\
        \hline
        Severe hypoglycemia & 0.8688\% \\
        Hypoglycemia & 3.0699\% \\
        Normoglycemia & 59.0382\% \\
        Hyperglycemia & 24.5741\% \\
        Severe hyperglycemia & 12.4490\% \\
        \hline
    \end{tabular}
\end{table}

To express the redistribution in absolute terms, Table~\ref{tab:normalized_baseline_adjustments} compares the original number of samples in each glycemic range with the number required to reach the normalized baseline distribution.

For dataset $d$ and glycemic range $r$, the target number of samples is defined as:

\[
n_{\mathrm{target}}(d,r)
=
N(d)\,p_{\mathrm{norm}}(r),
\]

and the required adjustment is:

\[
\Delta n(d,r)
=
n_{\mathrm{target}}(d,r)-n(d,r).
\]

A positive value indicates the number of samples that must be added through oversampling, whereas a negative value indicates the number of samples that must be removed through undersampling.

\begin{table}[htbp]
    \centering
    \caption{Absolute sample redistribution required to reach the normalized
    baseline distribution.}
    \label{tab:normalized_baseline_adjustments}
    \resizebox{\textwidth}{!}{%
    \begin{tabular}{llrrrl}
        \hline
        \textbf{Dataset} &
        \textbf{Glycemic range} &
        \textbf{Original samples} &
        \textbf{Target samples} &
        \textbf{Adjustment} &
        \textbf{Operation} \\
        \hline

        DiaTrend
        & Severe hypoglycemia
        & 35,291 & 66,611 & +31,320 & Oversampling \\

        DiaTrend
        & Hypoglycemia
        & 105,621 & 235,383 & +129,762 & Oversampling \\

        DiaTrend
        & Normoglycemia
        & 3,901,129 & 4,526,670 & +625,541 & Oversampling \\

        DiaTrend
        & Hyperglycemia
        & 2,192,705 & 1,884,184 & -308,521 & Undersampling \\

        DiaTrend
        & Severe hyperglycemia
        & 1,432,610 & 954,508 & -478,102 & Undersampling \\
        \hline

        REPLACE-BG
        & Severe hypoglycemia
        & 94,499 & 90,633 & -3,866 & Undersampling \\

        REPLACE-BG
        & Hypoglycemia
        & 291,856 & 320,268 & +28,412 & Oversampling \\

        REPLACE-BG
        & Normoglycemia
        & 6,636,660 & 6,159,107 & -477,553 & Undersampling \\

        REPLACE-BG
        & Hyperglycemia
        & 2,437,218 & 2,563,671 & +126,453 & Oversampling \\

        REPLACE-BG
        & Severe hyperglycemia
        & 972,175 & 1,298,729 & +326,554 & Oversampling \\
        \hline

        T1DiabetesGranada
        & Severe hypoglycemia
        & 220,610 & 193,155 & -27,455 & Undersampling \\

        T1DiabetesGranada
        & Hypoglycemia
        & 840,722 & 682,547 & -158,175 & Undersampling \\

        T1DiabetesGranada
        & Normoglycemia
        & 13,274,105 & 13,126,120 & -147,985 & Undersampling \\

        T1DiabetesGranada
        & Hyperglycemia
        & 5,281,558 & 5,463,626 & +182,068 & Oversampling \\

        T1DiabetesGranada
        & Severe hyperglycemia
        & 2,616,269 & 2,767,816 & +151,547 & Oversampling \\
        \hline
    \end{tabular}%
    }
\end{table}

These values represent the theoretical redistribution of the complete datasets. By applying this re-distribution all the three datasets will share common proportions amongst the glycemic ranges, regardless of their total size. There is some small residual differences caused by conversion to integer. These rounding differences will also be addressed in the implementation. 

\subsubsection{Fully Balanced Distribution}
\label{subsubsec:full_balancing_distribution}

Full balancing consists of forcing all glycemic ranges to contain the same number of samples, thereby completely removing the imbalance between the different glycemic regions. Accordingly, for the five glycemic ranges, the target proportion is defined as:

\[
p(r)=\frac{1}{5}=20\%.
\]

Thus, if (N) denotes the total number of training windows, each range is resampled to contain approximately (N/5) samples. Ranges below this target are oversampled, whereas ranges above it are undersampled. \textbf{\textit{This approach is considered the main strategy against which the study premise will be evaluated.}}

\subsubsection{Semi-Balanced Distribution}
\label{subsubsec:semi_balancing_distribution}

Semi-balancing is defined as an intermediate approach between the normalized baseline distribution and the fully balanced distribution. The aim of this strategy is to investigate the effects of reducing the imbalance without applying the more aggressive redistribution required by full balancing. However, this intermediate point is computed using a geometric mean rather than an arithmetic mean, corresponding to a midpoint on a logarithmic scale between the normalized baseline and the fully balanced distribution. The reason is that the objective is not to move each proportion halfway in absolute terms, but to reduce imbalance ratios halfway in relative terms.

Let \(p_{\mathrm{norm}}(r)\) denote the normalized baseline proportion of glycemic range \(r\). Since the fully balanced proportion is:

\[
p_{\mathrm{full}}(r)=\frac{1}{5}=0.20,
\]

The unnormalized semi-balanced value for each range is then defined as:

\[
q_{\mathrm{semi}}(r)
=
\sqrt{
p_{\mathrm{norm}}(r)
\cdot
p_{\mathrm{full}}(r)
}
=
\sqrt{
p_{\mathrm{norm}}(r)
\cdot 0.20
}.
\]

These values are then normalized so that their sum is equal to one:

\[
p_{\mathrm{semi}}(r)
=
\frac{
q_{\mathrm{semi}}(r)
}{
\displaystyle\sum_{k \in \mathcal{R}} q_{\mathrm{semi}}(k)
},
\]

where \(\mathcal{R}\) represents the set of five glycemic ranges.

As a result, applying these semi-balancing target proportions reduces the original imbalance ratios halfway on a logarithmic scale, while avoiding the stronger redistribution introduced by full balancing. The specific targets are shown in Table~\ref{tab:semi_balanced_target_proportions}

\begin{table}[htbp]
    \centering
    \caption{Semi-balanced target proportions.}
    \label{tab:semi_balanced_target_proportions}
    \begin{tabular}{lr}
        \hline
        \textbf{Glycemic range} & \textbf{Target proportion} \\
        \hline
        Severe hypoglycemia & 4.9438\% \\
        Hypoglycemia & 9.2934\% \\
        Normoglycemia & 40.7547\% \\
        Hyperglycemia & 26.2936\% \\
        Severe hyperglycemia & 18.7145\% \\
        \hline
    \end{tabular}
\end{table}

This approach has been considered because the literature also warns against applying balancing techniques too aggressively, especially when synthetic oversampling is used. Tarawneh et al. \cite{9761871} argue that, in particularly sensitive domains like this one, oversampling methods may generate synthetic examples that do not truly represent the intended minority region, leading to poor real-world behaviour. Similarly, Alkhawaldeh et al. \cite{alkhawaldeh2023challenges} also highlights that synthetic oversampling in medical datasets may introduce artificial instances that do not accurately represent the minority class, potentially producing unreliable results in real-world medical scenarios.

For these reasons, the three different strategies are experimentally explored: the normalized baseline provides a common sample-weighted EDA reference distribution, full balancing provides a naive upper reference where imbalance is completely removed, and semi-balancing provides an intermediate alternative.

\subsection{Balancing Methods}
\label{subsec:experimental_methodology_methods}

The main balancing process will be performed using SMOTER \cite{10.1007/978-3-642-40669-0_33} and SMOGN \cite{pmlr-v74-branco17a} because they are among the most established data-level methods specifically proposed for imbalanced regression problems, as stated in Chapter~\ref{chap:SoA}. While other imbalance-aware regression approaches exist, such as SERA \cite{5178874}, Balanced MSE \cite{Ren_2022_CVPR}, or geometric variants \cite{CAMACHO2022116387}, they either focus mainly on modifying the evaluation or loss function, or are less directly aligned with the objective of this thesis, which is to generate alternative balanced training sets while keeping the predictive models unchanged.

In addition to SMOTER \cite{10.1007/978-3-642-40669-0_33} and SMOGN \cite{pmlr-v74-branco17a}, a random resampling baseline will also be included in order to study a more naive approach. This method applies the same normalized, full-balancing and semi-balancing targets defined previously, but reaches them only through random duplication of minority-range windows and random removal of majority-range windows. Therefore it does not generate new synthetic glucose windows.

\subsection{Patient-Aware Balancing}
\label{subsec:experimental_methodology_patient_aware}

An additional methodological decision concerns whether balancing should be performed globally or independently for each patient. Since the EDA showed that patients do not share identical TBR, TIR, and TAR proportions, each individual presents a distinct glycemic profile shaped by their own physiological dynamics. Forcing every patient to reach the same balanced distribution would remove part of this heterogeneity and could generate patient profiles that do not reflect the original data.

For this reason, the proposed methodology will not apply strict per-patient balancing. Instead, the target distribution will be defined globally for each training set as defined in Chapter~\ref{subsec:target_distributions} while the resampling operations used to approach that target will be constrained by patient-level information. In other words, the objective is to balance the training set as a whole, not to force every patient to contain the same proportion of each glycemic range.

In practical terms, for each glycemic range, the number of samples to be added or removed will first be computed at global training-set level. When oversampling is required, only patients who already contain a minimum required ammount of 2 samples in the corresponding range will be considered eligible. The required additional samples will then be distributed across these eligible patients, and in proportion to their original contribution to that range.

\section{Statistical Analysis and Evaluation}
\label{sec:statistical_analysis}

The results of the experimentation will be statistically analysed in order to evaluate the study premise on a rigorous mathematical basis.

\subsection{Primary Clinical Outcomes And Effect Measure}
\label{subsec:primary_clinical_outcome}

The general hypothesis states that \textbf{increasing the representation of rare glycemic ranges during training has a positive effect on predictive performance in clinically critical regions}. However, this formulation is too broad, since it englobes the full cohort of the problem. For the scope of the statistical analysis, the hypothesis is therefore narrowed to the two hypoglycemic ranges, TBR$_2$ and TBR$_1$.

TBR$_2$ represents severe hypoglycemia and remains particularly relevant because it actually is the most clinically critical glycemic range -- as discussed in the previous chapters. However, it also contains a very small proportion of the available observations. Therefore, TBR$_1$ is also included as part of the statistical analysis in order to determine whether the effect of balancing is consistently observed in the broader hypoglycemic region.

Formally, for both ranges, the selected clinical outcome is the proportion of predictions located in Clarke Error Grid Zones A and B, denoted as:

\begin{equation}
AB_{\mathrm{TBR_2}},
\end{equation}

and

\begin{equation}
AB_{\mathrm{TBR_1}}.
\end{equation}

Zone A contains predictions considered clinically accurate, whereas Zone B contains predictions that deviate from the reference glucose value but would not normally lead to an inappropriate clinical decision. Therefore, combined, A+B represent the percentage of predictions that remain \textbf{clinically acceptable}.

The reference baseline for all datasets is defined by the results of the corresponding \texttt{NORMALIZED BASELINE} experiment, since it uses a common targets distribution for the three of them. Then, the fully balanced configuration is compared with the corresponding normalized configuration independently for both TBR$_2$ and TBR$_1$.

For each hypoglycemic range \(r\), balancing method \(m\) and cross-validation fold \(f\), the effect
is defined as:

\begin{equation}
\Delta AB_{r}^{(m,f)}
=
AB_{r}^{(m,f,\mathrm{full})}
-
AB_{r}^{(m,f,\mathrm{normalized})},
\label{eq:delta_ab_hypo}
\end{equation}

where:

\begin{equation}
r \in \{\mathrm{TBR_2}, \mathrm{TBR_1}\},
\end{equation}

\begin{equation}
m
\in
\{
\mathrm{Random},
\mathrm{SMOTER},
\mathrm{SMOGN}
\},
\end{equation}

and:

\begin{equation}
f \in \{1,2,3,4,5\}.
\end{equation}

A positive value of \(\Delta AB_{\mathrm{r}}^{(m,f)}\) indicates that full balancing increases the proportion of clinically acceptable hypoglycemia predictions relative to the normalized baseline for the corresponding hypoglycemic range, method and fold. A negative value indicates the opposite.

The comparison is paired by balancing method and cross-validation fold: each fully balanced result is compared with its normalized counterpart obtained using the same method and fold.

The three method-specific differences are subsequently averaged within each fold to minimize the impact of the balancing method on the analysis:

\begin{equation}
\Delta AB_{\mathrm{r}}^{(f)}
=
\frac{1}{3}
\sum_{m=1}^{3}
\Delta AB_{\mathrm{r}}^{(m,f)}.
\label{eq:fold_mean_delta_ab_hypo}
\end{equation}

The resulting five fold-level differences constitute the five observations used in
the statistical test.

\subsection{Statistical Hypothesis And Test}
\label{subsec:statistical_hypotheses}

Let \(\mu_{\Delta,r}\) denote the mean paired difference in \(AB_r\) between the fully balanced and normalized configurations across the five cross-validation folds, after averaging the corresponding differences across the three balancing methods, where:

\begin{equation}
r \in \{\mathrm{TBR_2},\mathrm{TBR_1}\}.
\end{equation}

Then, the null hypothesis states that full balancing does not produce a positive mean improvement relative to the normalized baseline:

\begin{equation}
H_{0,r}:
\mu_{\Delta,r} \leq 0.
\label{eq:null_hypothesis_hypo}
\end{equation}

The alternative hypothesis states that full balancing produces a positive mean improvement in hypoglycemia predictions clinical acceptability:

\begin{equation}
H_{1,r}:
\mu_{\Delta,r} > 0.
\label{eq:alternative_hypothesis_hypo}
\end{equation}

Rejecting ($H_{0,r}$) with a ($p$)-value below the selected significance level would provide enough evidence for our research premise: that full balancing improves \(AB_r\) relative to the normalized baseline after averaging across the evaluated balancing methods. If proven, this statistical evidence will then have to be evaluated together with any possible deterioration in other clincal metrics, and with the actual magnitude of \(\mu_{\Delta,r}\). This is to determine whether the observed change is practically meaningful -- \textit{satisfying Equation~\ref{eq:alternative_hypothesis_hypo} does not necessarily imply that the observed improvement is clinically relevant enough.}

Each hypothesis is evaluated independently using a one-sided paired Student's \(t\)-test. The paired design is appropriate because every fully balanced result is compared with the normalized result obtained using the same dataset and experimental configuration.

Let $d_{f,r}$ denote the paired difference for hypoglycemic range $r$ in each of the five folds, as defined in Equation~\ref{eq:fold_mean_delta_ab_hypo}. The mean paired difference is then calculated as:

\begin{equation}
\overline{d}_r
=
\frac{1}{n}
\sum_{f=1}^{n} d_{f,r},
\label{eq:mean_paired_difference}
\end{equation}

where \(n=5\) represents the five cross-validation folds.

The standard deviation of the paired differences is calculated as:

\begin{equation}
s_{d,r}
=
\sqrt{
\frac{
\sum_{f=1}^{n}
\left(
d_{f,r}-\overline{d}_r
\right)^2
}{
n-1
}
}.
\label{eq:paired_difference_sd}
\end{equation}

The paired Student's \(t\)-statistic is then defined as:

\begin{equation}
t_r
=
\frac{
\overline{d}_r
}{
s_{d,r}/\sqrt{n}
}.
\label{eq:paired_t_statistic}
\end{equation}

Under the null hypothesis, this statistic follows a Student's \(t\)-distribution with:

\begin{equation}
n-1=4
\end{equation}

degrees of freedom.

The significance level is established as:

\begin{equation}
\alpha=0.01.
\end{equation}

The null hypothesis will be rejected when:

\begin{equation}
p<\alpha.
\end{equation}

The paired \(t\)-test assumes that the fold-level paired differences are approximately normally distributed. However, in this case, this assumption cannot be formally assesed, since only five fold-level differences are available for each dataset and hypoglycemic range.

Additionally, since the same primary hypothesis is tested independently for both TBR$_2$ and TBR$_1$ across DiaTrend, REPLACE-BG, and T1DiabetesGranada, the resulting \(p\)-values will also be studied using the Holm procedure \cite{enwiki:1352765388} in order to also assess the probability of obtaining false-positive conclusions across the six dataset-specific tests.

\subsection{Secondary Outcomes}
\label{subsec:secondary_statistical_outcomes}

The primary confirmatory analysis focuses exclusively on the differences in \(AB_{\mathrm{TBR_2}}\) and \(AB_{\mathrm{TBR_1}}\) between full balancing and the normalized baseline. However, these outcomes alone cannot describe the complete effect of modifying the training-data distribution.

The semi-balanced configuration is retained as an intermediate approach and will be presented descriptively.

Additionally, other regression and Clarke Error Grid metrics are also considered as secondary outcomes. These include:

\begin{itemize}
    \item global RMSE, MAE, and A+B;
    \item TBR numerical performance and additional clinical metrics not included in the confirmatory analysis;
    \item TIR numerical and clinical performance;
    \item TAR numerical and clinical performance; and
    \item the individual proportions of predictions located in Clarke Error Grid Zones.
\end{itemize}

These secondary outcomes will be used to determine the impact of the primary clinical improvements -- if proven -- on other numerical prediction error metrics and on the movement of observations from clinically dangerous into acceptable ones. They are not treated as additional confirmatory hypotheses. They will mainly be interpreted through their numerical changes relative to the normalized baseline, and their consistency across balancing methods and datasets.

\section{Implementation and Experimental Execution}
\label{sec:experimental_execution}

\subsection{Balancing Module And Pipeline Integration}
\label{subsec:balancing_implementation}

The original training, validation, prediction, and evaluation pipelines had already been developed by the research team leading this thesis. Therefore, the objective is not to create a new prediction framework from scratch, but to incrust a new data-balancing mechanism into the existing framework so that the effect of different training-data distributions can be evaluated using the same training pipeline.

For this reason, the main contribution developed in this thesis is the BALANCING module. This module receives the training subset of each fold and generates a new version according to the selected balancing method and target distribution. Three methods were implemented: Random resampling, SMOTER and SMOGN. Each method can be combined with the normalized, semi-balanced, and fully balanced target distribution configurations described in Section~\ref{subsec:target_distributions}

The balancing process operates equivalently on the five clinically defined glycemic ranges. For each range, the number of required samples is calculated from the selected target distribution. Ranges containing more samples than required are undersampled, whereas minority ranges are oversampled: Random resampling generates additional observations by duplicating existing windows, while SMOTER and SMOGN generate synthetic windows from similar observations.

Additionally, for the synthetic approach (SMOTER and SMOGN), sample generation is patient-aware. The number of synthetic observations assigned to each patient is proportional to the number of eligible real observations available for that patient, and new windows are created using observations belonging to the same patient and glycemic range, rather than combining data from different patients.

Moreover, the inherited training pipeline was refactored to integrate this module. In each cross-validation fold, balancing is applied exclusively to the training subset before model fitting. The validation and test subsets remain unchanged so that model selection and final evaluation are performed using the original data distribution. The pipeline was also adapted to automatically execute the combinations of dataset, balancing method, and balancing level, while maintaining separate output directories and experiment records for every configuration.

All experiments use the same predictive model and training configuration. The model is a Long Short-Term Memory network with 128 recurrent units followed by a dense output layer that predicts a single future glucose value. It receives eight previous CGM measurements, representing two hours of glucose history, and predicts glucose four samples ahead, corresponding to a one-hour prediction horizon. The model is trained using the RMSE loss function and the Adam optimizer. Early stopping and model checkpointing are also used to retain the model with the best validation performance and avoid unnecessarily extending the training process.

A five-fold cross-validation procedure is also used for every experimental configuration. Results are calculated separately for the different glycemic ranges and are later weightedly-aggregated across folds. Consequently, the adapted pipeline provides a common and reproducible framework in which the only major experimental variation is the balancing strategy -- method and target -- applied to the training data.

Specific implementation can be checked in Appendix~\ref{app:source_code}.

\subsection{Experimental Workflow}
\label{subsec:experimental_workflow}

For each of the three datasets, the three target distributions described in Section~\ref{subsec:target_distributions} are combined with the three balancing methods presented in Section~\ref{subsec:experimental_methodology_methods}. This results in nine different experimental configurations per dataset, and 27 configurations in total:

\[
3\ \text{datasets}
\times
3\ \text{target distributions}
\times
3\ \text{balancing methods}
=
27\ \text{configurations}.
\]

Each configuration is evaluated using five-fold cross-validation. Within each fold, the selected balancing strategy is applied exclusively to the training subset, while the validation and test subsets remain untouched. The same fixed training configuration is used in every experiment. The resulting predictions are evaluated statistically as defined in Section~\ref{sec:statistical_analysis}.

\subsection{Model Training}
\label{sec:model_training}

The complete experimental execution are run on the Calculus computing server. Training runs are submitted as batch jobs through \textbf{SLURM} \cite{slurm_documentation}. This server-based workflow is necessary because of the large number of glucose windows and the computational cost of training all 27 -- 3 datasets x 3 methods x 3 targets -- balancing configurations across the different folds.